\subsection{TENSOR PRODUCT OF DIRECT LIMITS}

Let \((\mathbf{A}_{\alpha}, \rho_{\beta\alpha})\) be a direct system of rings and \((\mathbf{E}_{\alpha}, f_{\beta\alpha})\) (resp. \((\mathbf{F}_{\alpha}, g_{\beta\alpha})\) ) be a direct system of right (resp. left) \(A_{\alpha}\) -modules. For \(\alpha \leqslant \beta\) , there is a Z-module homomorphism

\[
f _ {\beta \alpha} \otimes g _ {\beta \alpha}: \mathrm{E} _ {\alpha} \otimes_ {\mathrm{A} _ {\alpha}} \mathrm{F} _ {\alpha} \rightarrow (\mathrm{E} _ {\beta}) _ {[ \mathrm{A} _ {\alpha} ]} \otimes_ {\mathrm{A} _ {\alpha}} (\mathrm{F} _ {\beta}) _ {[ \mathrm{A} _ {\alpha} ]}
\]

and on the other hand there is a canonical Z-module homomorphism

\[
\left(\mathrm{E} _ {\beta}\right) _ {[ \mathrm{A} _ {\alpha} ]} \otimes_ {\mathrm{A} _ {\alpha}} \left(\mathrm{F} _ {\beta}\right) _ {[ \mathrm{A} _ {\alpha} ]} \rightarrow \mathrm{E} _ {\beta} \otimes_ {\mathrm{A} _ {\beta}} \mathrm{F} _ {\beta}
\]

corresponding to the ring homomorphism \(\rho_{\beta \alpha}\) (§ 3, no. 3, Proposition 2); whence by composition we obtain a Z-module homomorphism

\[
h _ {\beta \alpha}: \mathrm{E} _ {\alpha} \otimes_ {\mathrm{A} _ {\alpha}} \mathrm{F} _ {\alpha} \rightarrow \mathrm{E} _ {\beta} \otimes_ {\mathrm{A} _ {\beta}} \mathrm{F} _ {\beta}
\]

which maps the tensor product \(x_{\alpha} \otimes y_{\alpha}\) to \(f_{\beta\alpha}(x_{\alpha}) \otimes g_{\beta\alpha}(y_{\alpha})\) . Clearly

\[
\left(\mathrm{E} _ {\alpha} \otimes_ {\mathrm{A} _ {\alpha}} \mathrm{F} _ {\alpha}, h _ {\beta \alpha}\right)
\]

is a direct system of \(\mathbf{Z}\) -modules. Let \(\mathbf{A} = \varinjlim \mathbf{A}_{\alpha}, \mathbf{E} = \varinjlim \mathbf{E}_{\alpha}, \mathbf{F} = \varinjlim \mathbf{F}_{\alpha}\) and let \(\rho_{\alpha}: \mathrm{A}_{\alpha} \to \mathrm{A}, f_{\alpha}: \mathrm{E}_{\alpha} \to \mathrm{E}, g_{\alpha}: \mathrm{F}_{\alpha} \to \mathrm{F}\) be the canonical mappings. As above, a Z-linear mapping \(\pi_{\alpha}: \mathrm{E}_{\alpha} \otimes_{\mathrm{A}_{\alpha}} \mathrm{F}_{\alpha} \to \mathrm{E} \otimes_{\mathrm{A}} \mathrm{F}\) is defined, which maps the tensor product \(x_{\alpha} \otimes y_{\alpha}\) to \(f_{\alpha}(x_{\alpha}) \otimes g_{\alpha}(y_{\alpha})\) , and it is immediate that these mappings form a direct system. Thus we obtain a Z-linear mapping

\[
\pi = \varinjlim \pi_ {\alpha}: \varinjlim (\mathrm{E} _ {\alpha} \otimes_ {\mathrm{A} \alpha} \mathrm{F} _ {\alpha}) \rightarrow \mathrm{E} \otimes_ {\mathrm{A}} \mathrm{F}.\tag{7}
\]

PROPOSITION 7. The Z-linear mapping (7) is bijective.

We set \(\mathrm{P} = \varinjlim (\mathrm{E}_{\alpha}\otimes_{\mathbf{A}_{\alpha}}\mathrm{F}_{\alpha})\) and, for all \(\alpha \in \mathbf{I}\) , let \(h_\alpha : \mathrm{E}_\alpha \otimes_{\mathbf{A}_\alpha}\mathrm{F}_\alpha \to \mathrm{P}\) be the canonical mapping. On the other hand, for all \(\alpha \in \mathbf{I}\) , let

\[
t _ {\alpha}: \mathrm{E} _ {\alpha} \times \mathrm{F} _ {\alpha} \rightarrow \mathrm{E} _ {\alpha} \otimes_ {\mathrm{A} _ {\alpha}} \mathrm{F} _ {\alpha}
\]

be the canonical Z-bilinear mapping; for \(\alpha \leqslant \beta\) ,

\[
t _ {\beta} \left(f _ {\beta \alpha} \left(x _ {\alpha}\right), g _ {\beta \alpha} \left(y _ {\alpha}\right)\right) = f _ {\beta \alpha} \left(x _ {\alpha}\right) \otimes g _ {\beta \alpha} \left(y _ {\alpha}\right) = h _ {\beta \alpha} \left(t _ {\alpha} \left(x _ {\alpha}, y _ {\alpha}\right)\right)
\]

and hence \((t_{\alpha})\) is a direct system of mappings. Canonically identifying \(\lim (\mathbf{E}_{\alpha}\times \mathbf{F}_{\alpha})\) with \(\mathbf{E}\times \mathbf{F}\) (Set Theory, III, § 7, no. 7, Proposition 10), we derive a mapping \(t = \lim t_{\alpha}:\mathbf{E}\times \mathbf{F}\to \mathbf{P}\) such that

\[
t \left(f _ {\alpha} (x _ {\alpha}), g _ {\alpha} (y _ {\alpha})\right) = h _ {\alpha} \left(t _ {\alpha} \left(x _ {\alpha}, y _ {\alpha}\right)\right) = h _ {\alpha} \left(x _ {\alpha} \otimes y _ {\alpha}\right).
\]

Taking account of Set Theory, III, § 7, no. 5, Lemma 1, it is immediately seen that \(t\) is Z-bilinear; moreover, for \(x \in \mathbf{E}\) , \(y \in \mathbf{F}\) , \(\lambda \in \mathbf{A}\) , there exists \(\alpha \in I\) such that \(x = f_{\alpha}(x_{\alpha})\) , \(y = g_{\alpha}(y_{\alpha})\) , \(\lambda = \rho_{\alpha}(\lambda_{\alpha})\) with \(\lambda_{\alpha} \in A_{\alpha}\) , \(x_{\alpha} \in E_{\alpha}\) , \(y_{\alpha} \in F_{\alpha}\) (Set Theory, III, § 3, no. 7, Lemma 1); whence

\[
t (x \lambda , y) = h _ {\alpha} ((x _ {\alpha} \lambda_ {\alpha}) \otimes y _ {\alpha}) = h _ {\alpha} (x _ {\alpha} \otimes (\lambda_ {\alpha} y _ {\alpha})) = t (x, \lambda y).
\]

Hence there exists one and only one Z-linear mapping \(\pi':\mathbf{E} \otimes_{\mathbf{A}} \mathbf{F} \to \mathbf{P}\) such that \(\pi'(x \otimes y) = t(x, y)\) (§ 3, no. 1, Proposition 1). Moreover, by definition

\[
\pi^ {\prime} \left(\pi \left(h _ {\alpha} \left(x _ {\alpha} \otimes y _ {\alpha}\right)\right)\right) = \pi^ {\prime} \left(f _ {\alpha} \left(x _ {\alpha}\right) \otimes g _ {\alpha} \left(y _ {\alpha}\right)\right) = h _ {\alpha} \left(x _ {\alpha} \otimes y _ {\alpha}\right)
\]

\[
\pi (\pi^ {\prime} (f _ {x} (x _ {\alpha}) \otimes g _ {\alpha} (y _ {\alpha}))) = \pi (h _ {\alpha} (x _ {\alpha} \otimes y _ {\alpha})) = f _ {\alpha} (x _ {\alpha}) \otimes g _ {\alpha} (y _ {\alpha})
\]

and as the elements of the form \(f_{\alpha}(x_{\alpha}) \otimes g_{\alpha}(y_{\alpha})\) (resp. \(h_{\alpha}(x_{\alpha} \otimes y_{\alpha})\) ) generate the Z-module E \(\otimes_{\mathbf{A}}\) F (resp. P), \(\pi' \circ \pi\) and \(\pi \circ \pi'\) are the identity mappings.

Loosely speaking, Proposition 7 may be expressed by saying that tensor products commute with direct limits and usually the two sides of (7) are identified by means of the isomorphism \(\pi\) .

COROLLARY 1. Let \((\mathbf{E}_{\alpha}', f_{\beta \alpha}')\) (resp. \((\mathbf{F}_{\alpha}', g_{\alpha \beta}'))\) be another direct system of right (resp. left) \(\mathbf{A}_{\alpha}\) -modules; for all \(\alpha \in \mathbf{I}\) , let \(u_{\alpha}: \mathbf{E}_{\alpha} \to \mathbf{E}_{\alpha}'\) (resp. \(v_{\alpha}: \mathbf{F}_{\alpha} \to \mathbf{F}_{\alpha}'\) ) be an \(\mathbf{A}_{\alpha}\) -linear mapping such that \((u_{\alpha})\) (resp. \((v_{\alpha})\) ) is a direct system. Then \((u_{\alpha} \oplus v_{\alpha})\) is a direct system of \(\mathbf{Z}\) -linear mappings and the diagram

\[
\begin{array}{c c c} \varinjlim (\mathrm{E} _ {\alpha} \otimes_ {\mathrm{A} _ {\alpha}} \mathrm{F} _ {\alpha}) & \longrightarrow & (\varinjlim \mathrm{E} _ {\alpha}) \otimes_ {\mathrm{A}} (\varinjlim \mathrm{F} _ {\alpha}) \\ \varinjlim (u _ {\alpha} \otimes v _ {\alpha}) \Bigg \downarrow & & \Bigg \downarrow (\varinjlim u _ {\alpha}) \otimes (\varinjlim v _ {\alpha}) \\ \varinjlim (\mathrm{E} _ {\alpha} ^ {\prime} \otimes_ {\mathrm{A} _ {\alpha}} \mathrm{F} _ {\alpha} ^ {\prime}) & \longrightarrow & (\lim \mathrm{E} _ {\alpha} ^ {\prime}) \otimes_ {\mathrm{A}} (\lim \mathrm{F} _ {\alpha} ^ {\prime}) \end{array}\tag{8}
\]

is commutative.

The verification is immediate.

Let \((\mathbf{A}_{\alpha}^{\prime}, \rho_{\beta\alpha}^{\prime})\) be another direct system of rings and suppose that each \(E_{\alpha}\) is an \((\mathbf{A}_{\alpha}^{\prime}, \mathbf{A}_{\alpha})\) -bimodule, the \(f_{\beta\alpha}\) being \((\mathbf{A}_{\alpha}^{\prime}, \mathbf{A}_{\alpha})\) -linear for \(\alpha \leqslant \beta\) . Then if we write \(A^{\prime} = \varinjlim A_{\alpha}^{\prime}\) , the isomorphism (7) is linear with respect to the left \(A^{\prime}\) -module structures on the two sides by virtue of Corollary 1. This can be immediately generalized to arbitrary multimodules.

In particular, if the \(A_{\alpha}\) are commutative, \(A = \varinjlim A_{\alpha}\) is commutative and isomorphism (7) is an A-module isomorphism.

COROLLARY 2. Let \((\mathbf{E}_{\alpha}, f_{\beta \alpha})\) be a direct system of right \(\mathbf{A}_{\alpha}\) -modules and let \(\mathbf{E}_{\alpha}' = \mathbf{E}_{\alpha} \otimes_{\mathbf{A}_{\alpha}} \mathbf{A}\) be the A-module obtained by extending the ring of scalars to \(\mathbf{A} = \lim_{\mathbf{A}_{\alpha}} \mathbf{A}_{\alpha}\) by means of the canonical homomorphism \(\rho_{\alpha}: \mathbf{A}_{\alpha} \to \mathbf{A}\) . Then \((\mathbf{E}_{\alpha}', f_{\beta \alpha} \otimes 1_{\mathbf{A}})\) is a direct system of right A-modules, whose direct limit is canonically isomorphic to \(\lim_{\mathbf{E}_{\alpha}} \mathbf{E}_{\alpha}\) .

It suffices to apply Proposition 7 with \(F_{\alpha}\) the ring A considered as an \((\mathbf{A}_{\alpha}, \mathbf{A})\) -bimodule by means of \(\rho_{\alpha}\) .

COROLLARY 3. Let A be a ring, \((\mathrm{E}_{\alpha}, f_{\beta\alpha})\) a direct system of right A-modules and F a left A-module. Then the Z-modules \(\lim_{\longrightarrow} (\mathrm{E}_{\alpha} \otimes_{\mathrm{A}} \mathrm{F})\) and \((\lim_{\longrightarrow} \mathrm{E}_{\alpha}) \otimes_{\mathrm{A}} \mathrm{F}\) are canonically isomorphic.

It suffices to take \(A_{\alpha} = A\) and \(F_{\alpha} = F\) for all \(\alpha \in I\) in Proposition 7.

In particular, if \(\rho: \mathbf{A} \to \mathbf{B}\) is a ring homomorphism, \(\lim_{\rho^* (\mathbf{E}_\alpha)} \rho^*(\mathbf{E}_\alpha)\) and \(\rho^*(\lim_{\rho \to \infty} \mathbf{E}_\alpha)\) are canonically isomorphic.

COROLLARY 4. Let M be a right A-module, N a left A-module, \((x_{i})_{1\leqslant i\leqslant n}\) a family of elements of M, \((y_{i})_{1\leqslant i\leqslant n}\) a family of elements of N, such that \(\sum_{i}(x_{i}\otimes y_{i}) = 0\) in M \(\otimes_{\mathbf{A}}\mathbf{N}\) . Then there exists a finitely generated submodule \(\mathbf{M}_1\) (resp. \(\mathbf{N}_1\) ) of M (resp. N) containing the \(x_{i}\) (resp. the \(y_{i}\) ) and such that \(\sum_{i}(x_{i}\otimes y_{i}) = 0\) in \(\mathbf{M}_1\otimes_{\mathbf{A}}\mathbf{N}_1\) .

M (resp. N) is canonically identified with the direct limit of the right directed family of its finitely generated submodules containing the \(x_{i}\) (resp. the \(y_{i}\) ) and it suffices to apply Set Theory, III, § 7, no. 5, Lemma 1.

